\documentclass[12pt]{article}
\usepackage{graphicx} % Required for inserting images
\usepackage[a4paper,margin=2cm]{geometry}
\usepackage{hyperref}

\begin{document}
\begin{titlepage}
    \centering
    \vspace*{2cm}
    {\Huge \textbf{B201B Comupter Science Lab (SS0326)}}
\vspace{3cm}

    {\Large Report}
   
    {\Large Tony Gamiel Naim Kaliny Armal (GH1049917)}
    
    {\Large GISMA University of Applied Sciences}
    
    {\Large Lecturer: Dr. William Baker Morrison}
    
{\Large 25/6/2026}
\vspace{3cm}






\includegraphics[width=1\textwidth]{714-gisma_logo550x241.png}

\vfill    
\end{titlepage}
%%%%%%%%%%%%%%%%%%
\pagestyle{plain}
\pagenumbering{arabic}
%%%%%%%%%%%%
\tableofcontents
\newpage
%%%%%%%%%%
\section*{Introduction}
\addcontentsline{toc}{section}{Introduction}
This report explains the design and development of my online portfolio for this module. The man goal of the assesment is to create a personal portfolio website that has all information about me and the Latex CV in clear and clean way. Nowdays many companies and employeres use online portfolios to evaluate interviewers before coming to the interview. The protfolio allows to show the skills, projects, certificates, etc. And for this reason doink or creating a portfolio is so important for students that study computer related stuff. For this project I created portfolio website using HTML and for my CV and report I used LaTeX. As also Github pages was used to host and publish my portfolio and beside that I did a Github repository that contains all the files related to the project the website code, images and my LaTeX code for the report and the CV. The project also gave me so much skills and experience in lots of branches such as Github, web development and LaTeX documentation as this report shows the technologies used, my repository and CV.



%%%%%%%%%%%%%%%%%%%%%
\section*{Portfolio Overview}
\addcontentsline{toc}{section}{Portfolio Overview}
The portfolio website was created to provide simple introduction about my self. As it also contains my education, skills and contact details.

\vspace{0.4cm}


\newline
Portfolio Website Link: https://tonygamiel997-dev.github.io/Tony-Armal.github.io/
\vspace{0.4cm}

The home page intrduces me shortly as software Enineering Student and explains my intrests in one or two words also wanted the visitor to understand quick who am I and what am I studying. The about me section has short information of the description of my background and interests. Another important section is the skills section as this area shows my skills that I gained or I have as also I am still developing my skills.
The education section includes current statues about me and my high school also. The contact section allows the visitors to reach or contact me through many was stated there in the website as number, email, github, my CV and linked in. These ioptions make it easier to contact or reach me. Overall the website is designed as a simple way to scroll and navigate. As also the main goal is to make sure that the visitors can find what the need easily without searching alot and saving time for them. 


%%%%%%%%%%%%%%%%%%%%
\section*{Github Repository}
\addcontentsline{toc}{section}{Github Repository}
The portfolio project is stored in the github repository and also Github is one of the biggest platforms that are used in the software development and project management (GitHub, 2025). It allows developers to store, organize and share the project files.

\vspace{0.4cm}
GitHub Repository: https://github.com/TonyGamiel997-dev/Tony-Armal.github.io
\newline
The repository contains all the files related to the project and the report too. Includes the code, images, files, LaTeX CV and Report. The repository is organised in a way that allows the users or developers understand the structure of the project. One also of the important files is the README file that contains brief explanation of the project and the contents as also it is considerd as good practice because it helps visitors the repository before opening it or exploring the files. GitHub also allows the developers track changes or change and edit the work (GitHub, 2025). However this project is small but using or learning how to use GitHub is so important in the software development.
%%%%%%%%%%%%
\section*{Job Market}
\addcontentsline{toc}{section}{Job Market}
One of the main reasons also for creating this portfolio was the relevance to the job market and future software engineering careers. As many companies now look at portfolios and repositories when evaluating candidates. My portfolio includes liks to my GitHub and linked in profile as there is downloadable button for my CV also. Linkedln is used for job opportunities (Linkedln, 2025).

%%%%%%%%%%%%%%%%%%
\section*{Tools used}
\addcontentsline{toc}{section}{Tools used}
Several tools and technologies were used to develope this portfolio.
\begin{itemize}
    \item HTML ( to create the website)
    \item GitHub ( was used to host my page and publish my page online as it also contains my repository and files for the project and the report)
    \item LaTeX ( was used to create my Cv and this report)
    
\end{itemize}
 The mix of these technologies allowed me to make a porfolio, CV and report while also at the same time I was gaining practical experience.
 
%%%%%%%%%

\section*{Improvments}
\addcontentsline{toc}{section}{Improvments}
There are lots of improvments I would like to do or add in the future and the most important is adding mnore project taht I wil develop or create along my studies such as from assignments or personal projects I will include all. I will also add new programming languages and technologies that I will learn along my bachelor. Another improvement detailed project description s as this will provide visitors with better understanding to my technical abilities. I also plan to include certificates maybe internships and achievments and other experiences as I progress in my academic career. By updating my website everytime I will become a stronger developer and represent my skills and groth as Software Engineering Student.

%%%%%%%%%%%%
 \section*{Conclusion}
\addcontentsline{toc}{section}{Conclusion}

This project gave me good experience in web development, GitHub and LaTeX also. The website portfolio presented most of my information in clear way as also through the the use of HTML, GitHub pages and LaTeX as mentioned above in the previous parts. Overall the project helped me to improve my skills work more with codes and learn more about programming.
%%%%%%%
\newpage
\section*{References:}
\addcontentsline{toc}{section}{References}
GitHub (2025) GitHub: Build and ship software on a single platform. Available at:  https://github.com/ (Accessed: 25 June 2026).

\vspace{0.2cm}

GitHub Pages (2025) GitHub Pages Documentation. Available at:  https://pages.github.com/ (Accessed: 25 June 2026).

\vspace{0.2cm}

LinkedIn (2025) Linked In Professional Networking Platform. Available at:  https://www.linkedin.com/ (Accessed: 25 June 2026).






%%%%%%%%%%%%%%%






\end{document}
